% -*- coding: utf-8 -*-
% !TEX program = xelatex
\documentclass[plain,noanswer,medmath]{../../template/jnuexam}
\usepackage{../../template/kysx}

\begin{document}


\examtitle{2005}{2}


\loadproblems{1}{2005P1}%载入试卷一


\makepart{填空题}{1～6小题，每小题4分，共24分}


\begin{problem}
设$y = (1 + \sin x)^x$，则$\dy\big|_{x = \pi} =$ \fillin{}．
\end{problem}


\begin{problem}
曲线$y = \frac{(1 + x)^{3/2}}{\sqrt{x}}$的斜渐近线方程为 \fillin{}．
\end{problem}


\begin{problem}
$\int_0^1 \frac{x\dx}{(2 - x^2)\sqrt{1 - x^2}} =$ \fillin{}．
\end{problem}


\useproblem{1}{1}{2}%【同试卷一第一(2)题】


\begin{problem}
当$x \to 0$时，$\alpha (x) = kx^2$与$\beta (x) = \sqrt{1 + x\arcsin x} - \sqrt{\cos x}$是等价无穷小，
则$k =$ \fillin{}．
\end{problem}


\useproblem{1}{1}{5}%【同试卷一第一(5)题】


\makepart{选择题}{7～14小题，每小题4分，共32分}


\useproblem{1}{2}{7}%【同试卷一第二(7)题】


\useproblem{1}{2}{8}%【同试卷一第二(8)题】


\begin{problem}
设函数$y = y(x)$由参数方程$\begin{cases}
  x = t^2 + 2t, \\
  y = \ln(1 + t)
\end{cases}$确定，则曲线$y = y(x)$在$x = 3$处的法线与$x$轴交点的横坐标是\pickout{}
\begin{abcd}
\item $\frac{1}{8}\ln 2 + 3$．
\item $- \frac{1}{8}\ln 2 + 3$．
\item $- 8\ln 2 + 3$．
\item $8\ln 2 + 3$．
\end{abcd}
\end{problem}


\begin{problem}
设区域$D = \left\{(x,y) \middle| x^2 + y^2 \le 4,x \ge 0,y \ge 0 \right\}$，$f(x)$为$D$上的正值连续函数，
$a,b$为常数，则$\iint_D \frac{a\sqrt{f(x)} + b\sqrt{f(y)}}{\sqrt{f(x)} + \sqrt{f(y)}} \d\sigma=$\pickout{}
\begin{abcd}
\item $ab\pi$．
\item $\frac{ab}{2}\pi$．
\item $(a + b)\pi$．
\item $\frac{a + b}{2}\pi$．
\end{abcd}
\end{problem}


\useproblem{1}{2}{9}%【同试卷一第二(9)题】


\begin{problem}
设函数$f(x) = \frac{1}{\e^{\frac{x}{x - 1}} - 1}$，则\pickout{}
\begin{abcd}
\item $x = 0$, $x = 1$都是$f(x)$的第一类间断点．
\item $x = 0$, $x = 1$都是$f(x)$的第二类间断点．
\item $x = 0$是$f(x)$的第一类间断点，$x = 1$是$f(x)$的第二类间断点．
\item $x = 0$是$f(x)$的第二类间断点，$x = 1$是$f(x)$的第一类间断点．
\end{abcd}
\end{problem}


\useproblem{1}{2}{11}%【同试卷一第二(11)题】


\useproblem{1}{2}{12}%【同试卷一第二(12)题】


\makepart{解答题}{15～23小题，共94分}


\begin{problem}[points=11]
设函数$f(x)$连续，且$f(0) \ne 0$，
求极限$\lim_{x \to 0} \frac{\int_0^x (x - t)f(t)\dt}{x\int_0^x f(x - t)\dt}$．
\end{problem}


\begin{problem}[points=11]
如图，$C_1$和$C_2$分别是$y = \frac{1}{2}(1 + \e^x)$和$y = \e^x$的图象，
过点(0,1)的曲线$C_3$是一单调增函数的图象．
过$C_2$上任一点$M(x,y)$分别作垂直于$x$轴和$y$轴的直线$l_x$和$l_y$．
记$C_1$, $C_2$与$l_x$所围图形的面积为$S_1(x)$；$C_2$, $C_3$与$l_y$所围图形的面积为$S_2(y)$．
如果总有$S_1(x) = S_2(y)$，求曲线$C_3$的方程$x = \phi(y)$．
\par\centerline{\includegraphics[width=0.5\linewidth]{2005B1}}
\end{problem}


\useproblem[points=11]{1}{3}{17}%【同试卷一第三(17)题】


\begin{problem}[points=12]
用变量代换$x = \cos t$（$0 < t < \pi$）化简微分方程$(1 - x^2)y'' - xy' + y = 0$，
并求其满足$y\big|_{x = 0} = 1$, $y'\big|_{x = 0} = 2$的特解．
\end{problem}


\useproblem[points=12]{1}{3}{18}%【同试卷一第三(18)题】


\begin{problem}[points=10]
已知函数$z = f(x,y)$的全微分$\dz = 2x\dx - 2y\dy$，并且$f(1,1) = 2$． 
求$f(x,y)$在椭圆域$D = \left\{(x,y)\middle| x^2 + \frac{y^2}{4} \le 1 \right\}$上的最大值和最小值．
\end{problem}


\begin{problem}[points=9]
计算二重积分$\iint_D \left| x^2 + y^2 - 1 \right| \d\sigma$，
其中$D = \left\{(x,y)\middle| 0 \le x \le 1,0 \le y \le 1 \right\}$
\end{problem}


\begin{problem}[points=9]
确定常数$a$，使向量组$\alpha_1 = (1,1,a)^T$, $\alpha_2 = (1,a,1)^T$, $\alpha_3 = (a,1,1)^T$
可由向量组$\beta_1 = (1,1,a)^T$, $\beta_2 = (- 2,a,4)^T$, $\beta_3 = (- 2,a,a)^T$线性表示，
但向量组$\beta_1$, $\beta_2$, $\beta_3$不能由向量组$\alpha_1$, $\alpha_2$, $\alpha_3$线性表示．
\end{problem}


\useproblem[points=9]{1}{3}{21}%【同试卷一第三(21)题】


\end{document}
